\begin{corollary}[7.1.4]
\label{IV.2.s7.1.4-fr}
Sea $A$ un anillo local noetheriano formalmente
equidimensional. Entonces:
\begin{enumerate}
\item[\textup{(i)}] $A$ es equidimensional y catenario
\textup{(es decir,} biequidimensional\textup{)}.
\item[\textup{(ii)}] Sea $\mathfrak a$ un ideal de $A$.
Para que $A/\mathfrak a$ sea equidimensional es
necesario y suficiente que $A/\mathfrak a$ sea
formalmente equidimensional. En particular, para
todo ideal primo $\mathfrak p$ de $A$,
$A/\mathfrak p$ es formalmente equidimensional.
\end{enumerate}
\end{corollary}

\begin{proof}
Si $\mathfrak m$ es el ideal maximal de $A$,
$\mathfrak m\widehat A$ es un ideal de definición
de $\widehat A$. Por hipótesis, $\widehat A$ es
equidimensional, y se sabe además que es catenario
(5.6.4); basta, pues, aplicar (7.1.3) a
$A'=\widehat A$.
\end{proof}
